\begingroup
\makeatletter\tagsleft@true\let\veqno\@@leqno\makeatother
\ArithmeticHead{266}{L}{ANALYSIS. NACHLASS.}
\SourceLeaf{273b}
\pdfbookmark[1]{Erläuterungen zu [III.]}{x1-agm-erlaeuterung-iii}
\begin{center}\textls[120]{Erläuterungen zu [III.], S. 181—183.}\end{center}
\par\indent Die Reihe $\varphi x$ des ersten Zettels art. [1.] geht aus der Reihe (18) des \textit{Abrisses} (S. 254) hervor, indem man darin $k'=\frac1x$ setzt und dann noch mit $\frac1x$ multipliziert; sie genügt derselben Differentialgleichung (20) S. 254, der die Reihen (18), (19) Genüge leisten. Mit Hilfe der Differentialgleichung wird die Funktionalgleichung (1) des art. [2.], S. 182 hergeleitet, aus der dann im zweiten Zettel art. [2.] erschlossen werden soll, daß $\varphi$ den reziproken Wert eines agM. darstellt, so daß es sich also um eine neue Ableitung der Reihenentwicklung für $\frac1{M(1,k)}$ handelt. Es kommt dabei auf den Schluß an, daß eine Funktion $F(z)$, die der Funktionalgleichung
\begin{equation*}\tag*{$(\alpha)$}F\left(\frac{1+t^2}{2t}\right)=F(t^2)\end{equation*}
\noindent Genüge leistet, eine Konstante sein müsse (S. 183). Setzt man
\[\xi=\frac1t,\quad G\left(\frac1{t^2}\right)=F(t^2),\]
\noindent so lautet die Gleichung $(\alpha)$:
\clearpage
\fancyhead{}\fancyhead[R]{\small267}\fancyhead[C]{\scriptsize BEMERKUNGEN ZUM AGM. ABSCHNITT III.}
\SourceLeaf{274}
\begin{equation*}\tag*{$(\beta)$}G\left(\frac{2\xi}{1+\xi^2}\right)=G(\xi^2).\end{equation*}
\noindent Bildet man den Algorithmus (5) des \textit{Abrisses} (S. 252) und setzt
\[k_n=\frac{\bar c_n}{a_n}=\frac{2\sqrt{k_{n-1}}}{1+k_{n-1}},\qquad(n=1,2,\ldots)\]
\noindent so folgt aus $(\beta)$
\[G(k_0)=G(k_1)=\ldots=G(k_n).\]
\noindent Nun ist aber nach (6) des \textit{Abrisses} (S. 252)
\[\lim_{n\to\infty}k_n=1.\]
\noindent Wenn also die Funktion $G(x)$ so beschaffen ist, daß $\displaystyle\lim_{x\to1}G(x)$ als endlicher oder unendlich großer Wert existiert, der unabhängig ist von dem Wege, auf dem die Veränderliche $x$ in den Punkt $x=1$ einrückt, so ist für ein jedes $k_0$, für das $G(k_0)$ endlich und stetig ist,
\[G(k_0)=\lim_{n\to\infty}G(k_n)=\lim_{x\to1}G(x),\]
\noindent also in der Tat konstant. Dies ist aber nicht notwendig der Fall, wenn die Funktion $G(x)$ die angegebene Bedingung für $x=1$ nicht erfüllt; so befriedigt z. B. die Funktion
\[G(x)=\left\{\frac{iM(1,\sqrt{1-x^2})}{M(1,x)}\right\}^{\frac{2\pi i}{\log2}},\]
\noindent für die jene Bedingung nicht erfüllt ist, die Funktionalgleichung $(\beta)$, ohne sich auf eine Konstante zu reduzieren.
\par\indent Wir bemerken noch, daß die Funktionalgleichung für $M(t)=M(t,1)$, von der im art. [2.] Gebrauch gemacht wird, vollständig mit der übereinstimmt, die im art. 5. der Abhandlung von 1800 (Werke III, S. 365) benutzt wird.
\par\indent Aus der Feststellung, daß der reziproke Wert von $M(\surd1+y,1)$ mit dem von $V$ freien Gliede der Entwicklung von
\[\frac1{\sqrt{1+y\sin^2 V}}\]
\noindent übereinstimmt (siehe S. 183), folgt für $y=1$, daß der reziproke Wert von $M(\surd2,1)$ gleich dem absoluten Gliede der Entwicklung von $(1+\sin^2 V)^{-\frac12}$ sei. Von diesem absoluten Gliede hatte \textsc{Gauss} aber schon zu Anfang des Jahres 1797 bei S. 54 des \textsc{Leiste} eingeschrieben (siehe oben S. 149, art. [8.], daß es gleich
\[\frac2\pi\int_0^1\frac{dx}{\sqrt{1-x^4}}=\frac\varpi\pi\]
\noindent sei; vergl. auch die Aufzeichnung auf der Rückseite des Leistetitelblatts, oben S. 169, art. [5.]. Wir haben damit alles, was für einen Beweis der Bemerkung in der Nr. 98 des \textit{Tagebuchs} (siehe oben S. 260) erforderlich ist, hier nachgewiesen.
\par\indent Am Schluß des art. [2.], S. 183 bedeuten $s$ und $c$ den Sinus und Kosinus eines Winkels, die kleine Tabelle gibt für die in der ersten Spalte stehenden Grade die Werte von $M(1,\sin\alpha^\circ)$; eine ausführlichere und genauere solche Tabelle ist im Abschnitt [IV.] art. [6.], S. 189 aus der Scheda Ac abgedruckt, vergl. die Bemerkungen zu diesem Abschnitte weiter unten.
\begin{flushright}34\textsuperscript{*}\end{flushright}
\clearpage
\endgroup

